\documentclass{article}
\usepackage[UTF8]{ctex}
\usepackage{amsmath}
\usepackage{amssymb}
\usepackage{graphicx}
\usepackage{geometry}
\geometry{a4paper, margin=2cm}

\newcommand{\problemname}[1]{\begin{center}\Large\bfseries #1\end{center}}
\usepackage{enumitem}

\begin{document}
\problemname{课前任务：打字练习}

\subsection*{题目描述}
本节为课前任务打字练习（重构题）。在代码编辑器中依次完成以下输入：第 1 行输入小写字母 \texttt{abc}；第 2 行按住 Shift 键输入大写字母 \texttt{DEF}；第 3 行点亮大写锁定键（Caps Lock）后输入大写字母 \texttt{GHI}；第 4 行熄灭大写锁定键后输入小写字母 \texttt{jkl}；第 5 行输入数字 \texttt{1230}；第 6 行输入符号 \texttt{<+-*/%;>}。请编写一个程序，不读入任何输入，将上述六行内容依次输出。

\subsection*{输入格式}
无输入。

\subsection*{输出格式}
依次输出六行内容：\\
abc\\
DEF\\
GHI\\
jkl\\
1230\\
<+-*/%;>

\subsection*{样例输入}
\begin{verbatim}

\end{verbatim}

\subsection*{样例输出}
\begin{verbatim}
abc
DEF
GHI
jkl
1230
<+-*/%;>
\end{verbatim}


\end{document}
